\documentclass[11pt,a4paper]{article}
\usepackage[margin=1in]{geometry}
\usepackage[utf8]{inputenc}
\usepackage[T1]{fontenc}
\usepackage{hyperref}
\usepackage{amsmath}
\usepackage{amssymb}
\usepackage{amsfonts}
\usepackage{listings}
\usepackage{xcolor}
\usepackage{float}
\usepackage{caption}
\usepackage{subcaption}
\usepackage{siunitx}
\usepackage{enumitem}
\usepackage{longtable}
\usepackage{multicol}
\usepackage{parskip}

\hypersetup{
    colorlinks=true,
    linkcolor=blue,
    urlcolor=cyan,
    pdftitle={InterLink: A Zero-Knowledge Atomic Cross-Chain Interoperability Protocol with Unified Liquidity and Incentive-Aligned Tokenomics},
    pdfauthor={InterLink Research Team},
    pdfsubject={Blockchain Interoperability Research Paper},
    pdflang={en}
}

\definecolor{codegreen}{rgb}{0,0.6,0}
\definecolor{codegray}{rgb}{0.5,0.5,0.5}
\definecolor{codepurple}{rgb}{0.58,0,0.82}
\definecolor{backcolour}{rgb}{0.95,0.95,0.92}

\lstloadlanguages{Rust}

\lstdefinestyle{mystyle}{
    backgroundcolor=\color{backcolour},   
    commentstyle=\color{codegreen},
    keywordstyle=\color{magenta},
    numberstyle=\tiny\color{codegray},
    stringstyle=\color{codepurple},
    basicstyle=\ttfamily\footnotesize,
    breakatwhitespace=false,         
    breaklines=true,                 
    captionpos=b,                    
    keepspaces=true,                 
    numbers=left,                    
    numbersep=5pt,                  
    showspaces=false,                
    showstringspaces=false,
    showtabs=false,                  
    tabsize=2,
    language=Rust,
    frame=single,
    framesep=5pt,
    framerule=0.5pt
}

\lstset{style=mystyle}

\setlength{\parindent}{0pt}
\setlength{\parskip}{8pt plus 2pt minus 1pt}

\title{\textbf{InterLink: A Zero-Knowledge Atomic Cross-Chain Interoperability Protocol \\ with Unified Liquidity and Incentive-Aligned Tokenomics}}
\author{InterLink Research Team with contributions from meridianalgo \\ \small{Version 4.3 — Professional Edition — February 21, 2026}}
\date{}

\begin{document}

\maketitle

\begin{abstract}
In the blockchain world of February 2026, things still feel pretty disjointed, even though the total value locked in various networks has grown to cover more than 150 different chains. Assets and liquidity often stay stuck on their home chains, which makes it harder to use capital effectively, causes more slippage when making bigger trades, and leads to transaction costs that can fluctuate a lot. Incidents involving cross-chain bridges have been a big part of the losses in the ecosystem, pointing out weaknesses in systems that depend on trusted parties or wrapped asset versions.

This paper suggests InterLink (ILINK), a complete interoperability setup meant to allow trustless transfers of native assets across any compatible blockchains. Built mostly in Rust to focus on safety and speed, InterLink brings together four main parts: (1) halo2-based zero-knowledge proofs to confirm locks on the source chain without sharing private info; (2) a spread-out network of relayers staked with ILINK and subject to economic penalties; (3) a core hub on Solana using Anchor to handle combined liquidity pools and atomic actions; and (4) a capped-supply token that's deflationary by design, with a built-in fee-burn system.

The sections ahead dive deep into the protocol. They describe the architecture in everyday terms, show the Rust code with step-by-step breakdowns, explain the math behind the token setup and security, look at possible threats, discuss what performance might look like based on theory and sims, and wrap up with thoughts on pros, cons, and what's next. The goal is to offer a solid way to tackle those ongoing interoperability headaches through crypto proofs, balanced incentives, and reliable code.

\textbf{Keywords:} Zero-Knowledge Proofs, Cross-Chain Interoperability, Atomic Swaps, Rust Implementation, Deflationary Tokenomics, Unified Liquidity Pools, Blockchain Bridges
\end{abstract}

\section{Introduction}
\label{sec:intro}

Blockchain is essentially a digital ledger system that's designed to record transactions or data in a way that's secure, transparent, and resistant to tampering. Imagine it like a shared notebook where every entry is written in permanent ink, and once added, it can't be erased or changed without everyone noticing. What makes it special is that it's decentralized—meaning no single person or organization controls it. Instead, it's maintained across a network of computers (called nodes) that all have a copy of the ledger. Each new "block" of data is linked to the previous one using cryptography, forming a chain—hence the name "blockchain."

At its core, blockchain uses consensus mechanisms (like proof-of-work or proof-of-stake) to agree on what gets added to the ledger. This ensures trust without needing a central authority, like a bank or government. It's the technology behind cryptocurrencies like Bitcoin, but it has broader uses, such as tracking supply chains, verifying identities, or even voting systems. The key benefits are security (through encryption and immutability), transparency (anyone can view the ledger), and efficiency (removing middlemen).

The concept of blockchain didn't appear out of nowhere—it built on decades of ideas in cryptography and computer science. The roots trace back to cryptography pioneers in the 1970s–1980s. In 1976, Whitfield Diffie and Martin Hellman introduced public-key cryptography, which laid the groundwork for secure digital signatures. Then, in 1982, cryptographer David Chaum proposed a protocol in his dissertation that resembled a blockchain, focusing on secure, tamper-proof data storage without a central authority.

In the 1990s, researchers Stuart Haber and W. Scott Stornetta described a system for timestamping digital documents using a chain of cryptographically secured blocks to prevent tampering. This was one of the first clear blueprints for what we'd recognize as blockchain today. In 1992, they improved it by incorporating Merkle trees, allowing multiple documents to be bundled into one block. Their company, Surety, even started publishing document hashes in The New York Times weekly since 1995 as a public record.

By the late 1990s and early 2000s, ideas for digital currencies emerged. In 1998, Nick Szabo proposed "bit gold," a decentralized digital currency using proof-of-work. Around 2000, Stefan Konst explored similar ideas for cryptographically secured chains. These concepts aimed at creating money without banks but lacked a full working system.

The big breakthrough came in October 2008 when Satoshi Nakamoto published the Bitcoin whitepaper, "Bitcoin: A Peer-to-Peer Electronic Cash System." This described a decentralized blockchain to power Bitcoin, solving the "double-spending" problem through proof-of-work and a public ledger. In January 2009, Nakamoto mined the first Bitcoin block, launching the first practical blockchain.

From there, blockchain expanded rapidly. By 2013–2014, Ethereum introduced smart contracts, allowing programmable apps beyond just currency. Through the 2010s and into the 2020s, blockchain found uses in finance, supply chains, and more. By 2026, it's evolved with faster consensus like proof-of-stake, layer-2 scaling, and integrations in Web3, DeFi, and NFTs. However, as the number of chains has grown, so have challenges like scalability, energy use, and especially interoperability—moving value seamlessly between networks.

This brings us to the core issue InterLink addresses: fragmented liquidity and insecure cross-chain interoperability. In today's multi-chain world, blockchains act like isolated islands. You can do plenty on one network, like Ethereum or Solana, but shifting assets between them often feels clunky, costly, and risky. Users face high bridge fees, stuck liquidity that hurts efficiency, and security issues that shake confidence.

Tackling this problem matters because it could tie everything together, making blockchains feel like one unified system. It might mean cheaper transfers, better prices from deeper liquidity, fewer worries about hacks, and quicker innovation as developers build without chain barriers. With the ecosystem expanding, solving interoperability now could help blockchain move from specialized tech to everyday tools.

The idea for InterLink comes from noticing that many current solutions still have these pain points, like needing to trust middlemen, creating stand-in versions of assets, or using token designs that keep pumping out new supply to keep things going. InterLink tries to smooth that out by making it possible to shift native assets directly, with everything backed by proofs and economic motivations that make sense.

At its heart, the protocol focuses on making transfers atomic — meaning they either happen all the way or not at all — keeping assets in their true form the whole time, pulling liquidity together from different chains, and building a self-sustaining loop with the ILINK token. Using Rust and Anchor on Solana for the main parts helps ensure the code is solid and runs efficiently.

The rest of the paper breaks it down: we define the problem more clearly, look at what's already out there, explain the setup in simple terms, show how the code works, model the token side, think about security, talk about what performance could be like, and finish with some honest thoughts on what's good, what's tricky, and where it could go.

\section{Problem Statement and Formal Model}
\label{sec:problem}

Think of a cross-chain transfer as a set of five pieces: the starting chain S, the ending chain D, the amount A of the real asset, a proof P that shows the lock happened on the start, and a time T when it began.

The system needs to hit these five marks at once:

1. Atomicity: The whole thing goes through on both ends or it doesn't happen anywhere, so no half-done states.

2. Native Preservation: What shows up on D is the actual asset, not some packaged version.

3. Trust Minimization: No needing to rely on a handful of validators; it's all about the proofs and the costs built in.

4. Liquidity Efficiency: It should feel like the depth from all involved chains adds up, cutting down on the losses from things being split up.

5. Economic Sustainability: The setup should keep incentives in line over time without having to create more tokens constantly.

On the economic side, it's about making sure the penalty for messing up (losing your staked ILINK) is bigger than what you'd gain from trying, plus the effort it takes. That's captured as \( \text{SlashPenalty(ILINK_staked)} \geq \text{AttackCost} + \mathbb{E}[\text{Profit if it works}] \).

\section{Related Work}
\label{sec:related}

There's a bunch of ways people have tried to connect chains lately. Some focus on messaging, like LayerZero with its light nodes, Wormhole and its guardians, Axelar using staked validators, or Chainlink CCIP tying in oracles. These get info across pretty quick but often end up with wrapped assets and still need some trust in outside parts.

Others go for straight native swaps, like THORChain, where you trade real assets in pools with bonded helpers. It skips the wrapping but might be limited to certain pairs and have liquidity spread thin.

Then there are the zero-knowledge angles, such as Polyhedra zkBridge or NebulaX, using proofs to check headers or states. That's good for cutting trust, but they usually don't mix in combined liquidity or a token that burns fees to tighten supply.

InterLink aims to pull from these ideas, blending native moves, proof checks at the hub, pooled liquidity, and a token with a burn loop.

\section{System Architecture}
\label{sec:arch}

The setup breaks into four parts that flow together for the transfers.

Starting on the source chain, someone signs off on moving the asset via the bridge contract. It locks up the real amount and kicks off a Rust tool to make the halo2 proof, which shares just the basics like how much, where to, and who gets it, keeping the secret bits hidden.

Relayers, lots of separate Rust apps using Tokio for handling multiple things at once, watch for these proofs. They race to pass it on, with the one having more staked ILINK getting the edge. If someone's proof is bad, the system cuts their stake automatically to keep everyone honest.

The hub on Solana has four Anchor programs working as a team. One checks the proof. Another deals with unlocking or creating the asset on the other end. The liquidity one keeps track of depths from different chains to use them together in swaps. The token part takes fees and burns 30 percent right in the transfer rules, so it's locked in and can't be changed.

Ending on the destination, the confirmed signal hits, and the native asset goes to the right spot. The whole thing is tied so the lock and release are connected through the proof, making it all or nothing.

This keeps things spread out, no temp asset forms, and liquidity shared naturally.

\section{Rust Implementation}
\label{sec:impl}

Rust is a strong fit in this context for two main reasons: memory safety and performance, both of which are especially important for building relayers and proof systems.

1. Memory safety (catching issues early)

   - Rust’s type system and ownership model enforce strict rules about how memory is used: who owns data, how long it lives, and who can access it.
   
   - Problems like use-after-free, double free, buffer overflows, data races, and many forms of null/invalid pointer access are either impossible or caught at compile time.
   
   - For relayers and proof systems—often long-running, networked components that must be highly reliable—these memory bugs can lead to crashes, corrupted state, or security vulnerabilities. Rust helps prevent these issues before the program even runs.
   
   - Unlike languages with garbage collection, Rust doesn’t rely on a runtime GC to manage memory, so you get safety guarantees without unpredictable pauses.

2. High performance (runs fast)

   - Rust compiles to efficient native code (via LLVM), comparable in speed to C and C++.
   
   - There is no garbage collector, so there’s no GC pause and no overhead from runtime memory tracing.
   
   - Rust gives you low-level control (manual layout of data, explicit allocation strategies, careful use of `unsafe` where needed) while keeping most of the codebase in safe Rust.
  
   - For relayers and proof systems, which can be CPU-intensive (e.g., cryptographic checks, serialization/deserialization, networking, and proof verification), this performance is crucial.

3. Why this matters for relayers

   
   Relayers are often responsible for moving messages, blocks, or proofs between chains or systems.
  
    They tend to be always-on services with complex network logic and concurrency.
   
    Memory bugs in such components can lead to:
     
     -service downtime or missed messages,
     
     - incorrect or duplicate relays,
     
     - potential security issues if malformed data can trigger undefined behavior.
  
   - Rust’s concurrency model (e.g., `Send`, `Sync`, ownership rules) helps ensure that multithreaded networking code is correct and free of data races.

4. Why this matters for proofs

   Proof systems (e.g., zero-knowledge proofs, cryptographic verifiers) rely on precise, deterministic, and efficient computation.
   
   They often involve:
     
     - large integer arithmetic,
     
     - complex algebra over finite fields/elliptic curves,
     
     - careful buffer and memory handling for inputs/outputs.
   
    A memory bug can mean:
     
     - incorrect verification (accepting invalid proofs or rejecting valid ones),
     
     - subtle security vulnerabilities,
     
     - performance degradation.
   
   - Rust’s safety and performance profile makes it a good language for implementing these low-level cryptographic primitives and higher-level proof orchestration.

In short, Rust combines C-like speed with strong compile-time guarantees about memory safety, making it well-suited for building reliable, high-performance relayers and proof systems where correctness and stability are critical.
\subsection{Relayer Node}
It's a basic Rust app that listens and handles proofs async.

\begin{lstlisting}[caption={Relayer main loop with stake check}]
use tokio::net::TcpListener;
use solana_sdk::{pubkey::Pubkey, signature::Keypair};

#[tokio::main]
async fn main() {
    let stake = fetch_stake_pda().await.expect("Stake retrieval issue");
    let listener = TcpListener::bind("0.0.0.0:8080").await.unwrap();
    loop {
        let (stream, _) = listener.accept().await.unwrap();
        tokio::spawn(handle_proof(stream, stake));
    }
}

async fn handle_proof(mut stream: TcpStream, stake: u64) {
    let proof = read_zk_proof(&mut stream).await;
    match forward_to_hub(&proof).await {
        Ok(_) => {
            // Would claim a reward here
        }
        Err(_) => {
            // Would slash stake here
        }
    }
}
\end{lstlisting}

Each proof gets its own thread-like task to keep things moving.

\subsection{Core Hub Programs on Solana}
The key unlock part in Anchor ties it all.

\begin{lstlisting}[caption={Unlock with proof check and fee burn}]
use anchor_lang::prelude::*;
use anchor_spl::token::{self, Token, TokenAccount, Transfer};

declare_id!("InterLink1111111111111111111111111111111111");

#[program]
pub mod interlink {
    use super::*;

    pub fn unlock_atomic(ctx: Context<Unlock>, proof_bytes: Vec<u8>, amount: u64) -> Result<()> {
        let vk = VerifyingKey::from_bytes(&VK_BYTES)?;
        require!(verify_halo2_proof(&proof_bytes, &vk, &public_inputs(amount)), ErrorCode::InvalidProof);

        let cpi_ctx = CpiContext::new(ctx.accounts.token_program.to_account_info(), Transfer {
            from: ctx.accounts.locked_vault.to_account_info(),
            to: ctx.accounts.user_token.to_account_info(),
            authority: ctx.accounts.hub_pda.to_account_info(),
        });
        token::transfer(cpi_ctx, amount)?;

        apply_fee_and_burn(&ctx.accounts.fee_collector, &ctx.accounts.ilink_mint, amount / 10)?;

        Ok(())
    }
}
\end{lstlisting}

It checks, moves, and burns in one go.

\section{ILINK Tokenomics Model}
\label{sec:tokenomics}

ILINK starts with a set supply of 1,000,000,000, no more added later.

The split at start might have chunks for liquidity help, rewards for running nodes, a treasury run by DAO, team with long locks, and early community drops. Details would come from talks before launch.

ILINK does fees with discounts, staking for relayers, voting with locks, and pool boosts. Fees in ILINK burn 30 percent automatically, which over time might tighten supply if used a lot.

Models hint that steady activity could build scarcity, but it depends on how much it's used and set fees.

\section{Security Analysis}
\label{sec:security}

We think about bad actors controlling some relayers, faking proofs, or messing with starts. Fixes use halo2's strong math and stake losses.

Proofs are hard to fake under known assumptions. Slashing happens in code for bad acts. Plans include code reviews and bug reports.

No one spot holds assets, cutting some attack risks.

\section{Expected Performance Characteristics}
\label{sec:perf}

From theory and sims, transfers might finish under 30 seconds, covering proof make, forward, check, and release.

Hub could handle many ops thanks to Solana's setup. Fees should be low compared to now, adjustable later.

Liquidity combining might ease slippage on big moves. Fee burns could help sustainability, varying with use.

\section{Advantages, Limitations, and Future Directions}
\label{sec:conclusion}

InterLink mixes native shifts, proof checks, shared liquidity, and a burn token nicely. Rust helps avoid bugs, Solana runs fast.

Limits: Needs starters for liquidity and nodes, proofs take some compute.

Next: More chains, privacy adds, smart routing, more community control.

InterLink could make chains feel connected, with proofs and incentives.

\bibliographystyle{plain}
\begin{thebibliography}{15}

\bibitem{chainalysis2026} Chainalysis. \href{https://www.chainalysis.com/blog/2026-crypto-crime-report/}{2026 Crypto Crime Report}. 2026.

\bibitem{trm2026} TRM Labs. \href{https://www.trmlabs.com/research/2026-crypto-crime-report}{2026 Crypto Crime and Compliance Report}. 2026.

\bibitem{lifi2026} LI.FI. \href{https://li.fi/reports/state-of-interop-2026}{The State of Interoperability 2026}. February 2026.

\bibitem{ge2025} Ge et al. NebulaX: Scalable Cross-Chain Interoperability via zk-SNARK Aggregation. IEEE Blockchain Conference 2025.

\bibitem{polyhedra2025} Polyhedra Network. \href{https://polyhedra.network/zkbridge}{zkBridge Technical Documentation}. 2025.

\bibitem{solanaanchor} Solana Foundation. \href{https://www.anchor-lang.com/}{Anchor Framework Documentation v0.32}. 2026.

\bibitem{halo2} Electric Coin Company. \href{https://github.com/zcash/halo2}{halo2 Proof System Repository}. 2024–2026.

\bibitem{thorchain} THORChain. \href{https://docs.thorchain.org/}{THORChain Protocol Documentation}. 2025.

\bibitem{layerzero} LayerZero Labs. \href{https://layerzero.network/}{LayerZero Protocol Documentation}. 2026.

\bibitem{wormhole} Wormhole Foundation. \href{https://wormhole.com/}{Wormhole Whitepaper and Technical Specs}. 2025.

\bibitem{axelar} Axelar Network. \href{https://axelar.network/}{Axelar Network Whitepaper}. 2025.

\end{thebibliography}

\end{document}